\documentclass[11pt]{article}
\usepackage[a4paper,margin=16mm,headheight=14pt,footskip=10mm]{geometry}
\usepackage[T1]{fontenc}
\usepackage{lmodern,microtype,amsmath,amssymb,booktabs,adjustbox,xcolor}
\usepackage{quantikz}
\usepackage{fancyhdr,hyperref}
\hypersetup{colorlinks=true,linkcolor=blue,urlcolor=blue,pdftitle={Full circuit: C4-Helix SWAP gauging},pdfauthor={}}
\definecolor{ink}{HTML}{203044}
\definecolor{cat}{HTML}{986F14}
\definecolor{gauge}{HTML}{148E9E}
\definecolor{groupa}{HTML}{367AAF}
\definecolor{groupb}{HTML}{D47D28}
\definecolor{groupc}{HTML}{378957}
\definecolor{groupd}{HTML}{8863AF}
\definecolor{groupe}{HTML}{CA5D85}
\setlength{\parindent}{0pt}
\setlength{\parskip}{5pt}
\setlength{\abovedisplayskip}{7pt}
\setlength{\belowdisplayskip}{7pt}
\pagestyle{fancy}\fancyhf{}
\fancyhead[L]{\small\color{ink}C4-Helix / full circuit}
\fancyhead[R]{\small Companion to full\_procedure.md}
\fancyfoot[L]{\footnotesize Quantum controls: solid wire. Classical feedforward: double wire.}
\fancyfoot[R]{\thepage}
\renewcommand{\headrulewidth}{.3pt}
\newcommand{\subhead}[1]{\par\vspace{3mm}{\bfseries\color{ink}#1}\par}
\begin{document}

{\Large\bfseries\color{ink}Shared gadget: reset, cat preparation, and verification}\par\vspace{3mm}

\small
Read every circuit left to right. A box $|0\rangle$ means reset, discarding
the previous state; it is not a projection with postselection. A meter
measures $Z$. The sequence $H$ then meter measures $X$.
Single wires are quantum; double wires carry classical bits.
\[
\mathcal R_0(\rho)=|0\rangle\langle0|\,\operatorname{Tr}\rho,
\qquad b=0\ \Longleftrightarrow\ \hbox{reported sign }+1.
\]
\subhead{Prepare a raw cat: the four-leaf example}

\begin{center}\begin{adjustbox}{max width=\linewidth}
\begin{quantikz}[row sep={.65cm,between origins},column sep=.55cm]
\lstick{$\color{cat}\ell_0$} & \gate{\lvert0\rangle} & \gate{H} & \ctrl{1} & \qw & \qw & \qw \\
\lstick{$\color{cat}\ell_1$} & \gate{\lvert0\rangle} & \qw & \targ{} & \ctrl{1} & \qw & \qw \\
\lstick{$\color{cat}\ell_2$} & \gate{\lvert0\rangle} & \qw & \qw & \targ{} & \ctrl{1} & \qw \\
\lstick{$\color{cat}\ell_3$} & \gate{\lvert0\rangle} & \qw & \qw & \qw & \targ{} & \qw
\end{quantikz}
\end{adjustbox}\end{center}

For $m$ leaves, extend the chain through $\ell_{m-1}$:
$H_{\ell_0}$ followed by CNOTs $\ell_j\to\ell_{j+1}$, $j=0,\ldots,m-2$.
\subhead{One adjacent-$ZZ$ verification, with its own reset}

\begin{center}\begin{adjustbox}{max width=\linewidth}
\begin{quantikz}[row sep={.6cm,between origins},column sep=.6cm]
\lstick{$\color{cat}\ell_j$} & \qw & \ctrl{2} & \qw & \qw & \qw \\
\lstick{$\color{cat}\ell_{j+1}$} & \qw & \qw & \ctrl{1} & \qw & \qw \\
\lstick{$\color{cat}w$} & \gate{\lvert0\rangle} & \targ{} & \targ{} & \meter{} & \rstick{$\nu_j^{(k)}$}\cw
\end{quantikz}
\end{adjustbox}\end{center}

\[
\text{For }k=1,2,3:\quad j=0,1,\ldots,m-2;\qquad
\text{require every }\nu_j^{(k)}=0.
\]
Reset $w$ before every pair. After verification, apply no further
leaf-to-leaf gates. Couple each leaf to one target port only, then measure
all leaves in $X$. The check bit is $y=\bigoplus_j\lambda_j$.
For a positive check require $y=0$, not every $\lambda_j=0$.

\newpage

{\Large\bfseries\color{ink}1. Prepare the encoded B and C inputs}\par\vspace{3mm}


\begin{center}\begin{adjustbox}{max width=\linewidth}
\begin{quantikz}[row sep={.6cm,between origins},column sep=.6cm]
\lstick{$\color{groupa}B_q$} & \gate{\lvert0\rangle} & \rstick{$q=0,\ldots,19$}\qw \\
\lstick{$\color{groupa}C_q$} & \gate{\lvert0\rangle} & \rstick{$q=0,\ldots,19$}\qw
\end{quantikz}
\end{adjustbox}\end{center}

The drawing represents an independent reset of every physical data qubit.
For each check below, prepare one fresh verified cat as on page 1.
The leaves in each panel are one joint cat, not independent plus states.
\par\noindent\begin{minipage}[t]{.485\linewidth}\centering\textbf{$X_{D_0}X_{D_1}X_{D_2}X_{D_3}$}

\begin{center}\begin{adjustbox}{max width=\linewidth}
\begin{quantikz}[row sep={.46cm,between origins},column sep=.22cm]
\lstick{$\color{cat}\ell_{0}$} & \ctrl{1} & \qw & \qw & \qw & \gate{H} & \meter{} & \rstick{$\lambda_{0}$}\cw \\
\lstick{$\color{groupa}D_0$} & \targ{} & \qw & \qw & \qw & \qw & \qw & \qw \\
\lstick{$\color{cat}\ell_{1}$} & \qw & \ctrl{1} & \qw & \qw & \gate{H} & \meter{} & \rstick{$\lambda_{1}$}\cw \\
\lstick{$\color{groupa}D_1$} & \qw & \targ{} & \qw & \qw & \qw & \qw & \qw \\
\lstick{$\color{cat}\ell_{2}$} & \qw & \qw & \ctrl{1} & \qw & \gate{H} & \meter{} & \rstick{$\lambda_{2}$}\cw \\
\lstick{$\color{groupa}D_2$} & \qw & \qw & \targ{} & \qw & \qw & \qw & \qw \\
\lstick{$\color{cat}\ell_{3}$} & \qw & \qw & \qw & \ctrl{1} & \gate{H} & \meter{} & \rstick{$\lambda_{3}$}\cw \\
\lstick{$\color{groupa}D_3$} & \qw & \qw & \qw & \targ{} & \qw & \qw & \qw
\end{quantikz}
\end{adjustbox}\end{center}
\end{minipage}\hfill\begin{minipage}[t]{.485\linewidth}\centering\textbf{$Z_{D_0}Z_{D_1}Z_{D_2}Z_{D_3}$}

\begin{center}\begin{adjustbox}{max width=\linewidth}
\begin{quantikz}[row sep={.46cm,between origins},column sep=.22cm]
\lstick{$\color{cat}\ell_{0}$} & \ctrl{1} & \qw & \qw & \qw & \gate{H} & \meter{} & \rstick{$\lambda_{0}$}\cw \\
\lstick{$\color{groupa}D_0$} & \control{} & \qw & \qw & \qw & \qw & \qw & \qw \\
\lstick{$\color{cat}\ell_{1}$} & \qw & \ctrl{1} & \qw & \qw & \gate{H} & \meter{} & \rstick{$\lambda_{1}$}\cw \\
\lstick{$\color{groupa}D_1$} & \qw & \control{} & \qw & \qw & \qw & \qw & \qw \\
\lstick{$\color{cat}\ell_{2}$} & \qw & \qw & \ctrl{1} & \qw & \gate{H} & \meter{} & \rstick{$\lambda_{2}$}\cw \\
\lstick{$\color{groupa}D_2$} & \qw & \qw & \control{} & \qw & \qw & \qw & \qw \\
\lstick{$\color{cat}\ell_{3}$} & \qw & \qw & \qw & \ctrl{1} & \gate{H} & \meter{} & \rstick{$\lambda_{3}$}\cw \\
\lstick{$\color{groupa}D_3$} & \qw & \qw & \qw & \control{} & \qw & \qw & \qw
\end{quantikz}
\end{adjustbox}\end{center}
\end{minipage}\par

\subhead{Full target lists and repetitions}
For $s=0,\ldots,8$, replace $(0,1,2,3)$ by the ordered support $R_s$
from page 13. An upper check has eight leaves. Logical checks have six.
\[
\begin{aligned}
\mathcal F_B={}&(S^X_{B,0:8},\ S^Z_{B,0:8},\ \overline Z^B_0,\overline Z^B_1),\\
\mathcal F_C={}&(S^X_{C,0:8},\ S^Z_{C,0:8},\ \overline X^C_0,\overline Z^C_1).
\end{aligned}
\]
The notation $0{:}8$ means all nine checks in increasing $s$ order.
Execute three complete positive rounds on B, then three on C.
Require all check bits zero, in addition to all cat-verification bits.
\[
|\psi\rangle=|\overline{00}\rangle_B|\overline{+0}\rangle_C.
\]

\newpage

{\Large\bfseries\color{ink}2. Prepare the graph resource and verify its edge checks}\par\vspace{3mm}

Reset $a,b,c,d,e$ and $A_0,\ldots,A_{19}$ to zero. Apply $H$
to each vertex. For every edge $q:u_q-v_q$, apply
CNOT$(u_q\to A_q)$ then CNOT$(v_q\to A_q)$, in increasing $q$.
\subhead{The gates for edges 0 and 1 (other edges continue the same wires)}

\begin{center}\begin{adjustbox}{max width=\linewidth}
\begin{quantikz}[row sep={.65cm,between origins},column sep=.43cm]
\lstick{$\color{groupa}a$} & \gate{\lvert0\rangle} & \gate{H} & \ctrl{3} & \qw & \qw & \qw & \qw \\
\lstick{$\color{groupc}c$} & \gate{\lvert0\rangle} & \gate{H} & \qw & \qw & \ctrl{3} & \qw & \qw \\
\lstick{$\color{groupe}e$} & \gate{\lvert0\rangle} & \gate{H} & \qw & \ctrl{1} & \qw & \ctrl{2} & \qw \\
\lstick{$\color{gauge}A_0$} & \gate{\lvert0\rangle} & \qw & \targ{} & \targ{} & \qw & \qw & \qw \\
\lstick{$\color{gauge}A_1$} & \gate{\lvert0\rangle} & \qw & \qw & \qw & \targ{} & \targ{} & \qw
\end{quantikz}
\end{adjustbox}\end{center}

\begin{center}\begin{tabular}{cc@{\qquad}cc}
Endpoints & Labels & Endpoints & Labels\\\hline
ab&7,15&bc&8,19\\ ac&11,12&bd&14,16\\
ad&3,6&be&4,9\\ ae&0,10&cd&2,13\\ ce&1,18&de&5,17
\end{tabular}\end{center}
\subhead{One edge-check measurement: $K_1=Z_cZ_{A_1}Z_e$}

\begin{center}\begin{adjustbox}{max width=\linewidth}
\begin{quantikz}[row sep={.46cm,between origins},column sep=.22cm]
\lstick{$\color{cat}\ell_{0}$} & \ctrl{1} & \qw & \qw & \gate{H} & \meter{} & \rstick{$\lambda_{0}$}\cw \\
\lstick{$\color{groupc}c$} & \control{} & \qw & \qw & \qw & \qw & \qw \\
\lstick{$\color{cat}\ell_{1}$} & \qw & \ctrl{1} & \qw & \gate{H} & \meter{} & \rstick{$\lambda_{1}$}\cw \\
\lstick{$\color{gauge}A_1$} & \qw & \control{} & \qw & \qw & \qw & \qw \\
\lstick{$\color{cat}\ell_{2}$} & \qw & \qw & \ctrl{1} & \gate{H} & \meter{} & \rstick{$\lambda_{2}$}\cw \\
\lstick{$\color{groupe}e$} & \qw & \qw & \control{} & \qw & \qw & \qw
\end{quantikz}
\end{adjustbox}\end{center}

One three-leaf verified cat is shared by these three ports.
Set $y=\lambda_0\oplus\lambda_1\oplus\lambda_2$ and require zero.
For edge $q:u-v$, use targets $(u,A_q,v)$.
The vertex checks on the next page complete resource verification.

\newpage

{\Large\bfseries\color{ink}2 continued. Verify a resource vertex check}\par\vspace{3mm}


\[
K_a=X_aX_{A_0}X_{A_3}X_{A_6}X_{A_7}
X_{A_{10}}X_{A_{11}}X_{A_{12}}X_{A_{15}}.
\]
Prepare \textbf{one nine-leaf cat}. The two panels show its first five
and last four target ports. They do not prepare or measure two separate cats.
\par\noindent\begin{minipage}[t]{.485\linewidth}\centering\textbf{Leaves 0--4}

\begin{center}\begin{adjustbox}{max width=\linewidth}
\begin{quantikz}[row sep={.46cm,between origins},column sep=.22cm]
\lstick{$\color{cat}\ell_{0}$} & \ctrl{1} & \qw & \qw & \qw & \qw & \gate{H} & \meter{} & \rstick{$\lambda_{0}$}\cw \\
\lstick{$\color{groupa}a$} & \targ{} & \qw & \qw & \qw & \qw & \qw & \qw & \qw \\
\lstick{$\color{cat}\ell_{1}$} & \qw & \ctrl{1} & \qw & \qw & \qw & \gate{H} & \meter{} & \rstick{$\lambda_{1}$}\cw \\
\lstick{$\color{gauge}A_{0}$} & \qw & \targ{} & \qw & \qw & \qw & \qw & \qw & \qw \\
\lstick{$\color{cat}\ell_{2}$} & \qw & \qw & \ctrl{1} & \qw & \qw & \gate{H} & \meter{} & \rstick{$\lambda_{2}$}\cw \\
\lstick{$\color{gauge}A_{3}$} & \qw & \qw & \targ{} & \qw & \qw & \qw & \qw & \qw \\
\lstick{$\color{cat}\ell_{3}$} & \qw & \qw & \qw & \ctrl{1} & \qw & \gate{H} & \meter{} & \rstick{$\lambda_{3}$}\cw \\
\lstick{$\color{gauge}A_{6}$} & \qw & \qw & \qw & \targ{} & \qw & \qw & \qw & \qw \\
\lstick{$\color{cat}\ell_{4}$} & \qw & \qw & \qw & \qw & \ctrl{1} & \gate{H} & \meter{} & \rstick{$\lambda_{4}$}\cw \\
\lstick{$\color{gauge}A_{7}$} & \qw & \qw & \qw & \qw & \targ{} & \qw & \qw & \qw
\end{quantikz}
\end{adjustbox}\end{center}
\end{minipage}\hfill\begin{minipage}[t]{.485\linewidth}\centering\textbf{Leaves 5--8}

\begin{center}\begin{adjustbox}{max width=\linewidth}
\begin{quantikz}[row sep={.46cm,between origins},column sep=.22cm]
\lstick{$\color{cat}\ell_{5}$} & \ctrl{1} & \qw & \qw & \qw & \gate{H} & \meter{} & \rstick{$\lambda_{5}$}\cw \\
\lstick{$\color{gauge}A_{10}$} & \targ{} & \qw & \qw & \qw & \qw & \qw & \qw \\
\lstick{$\color{cat}\ell_{6}$} & \qw & \ctrl{1} & \qw & \qw & \gate{H} & \meter{} & \rstick{$\lambda_{6}$}\cw \\
\lstick{$\color{gauge}A_{11}$} & \qw & \targ{} & \qw & \qw & \qw & \qw & \qw \\
\lstick{$\color{cat}\ell_{7}$} & \qw & \qw & \ctrl{1} & \qw & \gate{H} & \meter{} & \rstick{$\lambda_{7}$}\cw \\
\lstick{$\color{gauge}A_{12}$} & \qw & \qw & \targ{} & \qw & \qw & \qw & \qw \\
\lstick{$\color{cat}\ell_{8}$} & \qw & \qw & \qw & \ctrl{1} & \gate{H} & \meter{} & \rstick{$\lambda_{8}$}\cw \\
\lstick{$\color{gauge}A_{15}$} & \qw & \qw & \qw & \targ{} & \qw & \qw & \qw
\end{quantikz}
\end{adjustbox}\end{center}
\end{minipage}\par

\[
y(K_a)=\lambda_0\oplus\cdots\oplus\lambda_8;\qquad \text{require }y(K_a)=0.
\]
Use the same circuit for each vertex v, replacing a by v and
$E_a=R_5$ by $E_v$:
\[
K_v=X_v\prod_{q\in E_v}X_{A_q},\qquad
(E_a,E_b,E_c,E_d,E_e)=(R_5,R_6,R_7,R_8,R_9).
\]
\subhead{The complete resource-verification schedule}
Measure $(K_a,K_b,K_c,K_d,K_e,K_0,\ldots,K_{19})$ in that order,
with a fresh verified cat for every check. Execute three positive rounds.
These are $3\times25$ check results, all required zero.
The resource is $2^{-5/2}\sum_t|t\rangle_V|\delta t\rangle_A$.

\newpage

{\Large\bfseries\color{ink}3. Enrich the symmetry and measure the temporary vertices}\par\vspace{3mm}

The a vertex coherently controls four physical SWAPs. Each column
with one dot and two crosses is a three-qubit Fredkin.

\begin{center}\begin{adjustbox}{max width=\linewidth}
\begin{quantikz}[row sep={.58cm,between origins},column sep=.55cm]
\lstick{$\color{groupa}a$} & \ctrl{1} & \ctrl{3} & \ctrl{5} & \ctrl{7} & \gate{H} & \meter{} & \rstick{$b_a$}\cw \\
\lstick{$\color{groupa}B_0$} & \swap{1} & \qw & \qw & \qw & \qw & \qw & \qw \\
\lstick{$\color{groupa}C_0$} & \targX{} & \qw & \qw & \qw & \qw & \qw & \qw \\
\lstick{$\color{groupa}B_1$} & \qw & \swap{1} & \qw & \qw & \qw & \qw & \qw \\
\lstick{$\color{groupa}C_1$} & \qw & \targX{} & \qw & \qw & \qw & \qw & \qw \\
\lstick{$\color{groupa}B_2$} & \qw & \qw & \swap{1} & \qw & \qw & \qw & \qw \\
\lstick{$\color{groupa}C_2$} & \qw & \qw & \targX{} & \qw & \qw & \qw & \qw \\
\lstick{$\color{groupa}B_3$} & \qw & \qw & \qw & \swap{1} & \qw & \qw & \qw \\
\lstick{$\color{groupa}C_3$} & \qw & \qw & \qw & \targX{} & \qw & \qw & \qw
\end{quantikz}
\end{adjustbox}\end{center}

\begin{center}\begin{tabular}{ccl}
Vertex & Data-pair labels $J_v$ & Quantum action before its $X$ readout\\\hline
a&0--3&$\mathrm C_a(T_a)$\\
b&4--7&$\mathrm C_b(T_b)$\\
c&8--11&$\mathrm C_c(T_c)$\\
d&12--15&$\mathrm C_d(T_d)$\\
e&16--19&$\mathrm C_e(T_e)$
\end{tabular}\end{center}
Apply \textbf{all twenty Fredkins first}, increasing q. Then read vertices
a,b,c,d,e in X. The meter in the panel shows the a readout; defer it until
all twenty enrichment Fredkins have finished. The gauge qubits are idle
during enrichment.
\[
r_{\rm bit}=b_a\oplus b_b\oplus b_c\oplus b_d\oplus b_e;
\qquad \text{require }r_{\rm bit}=0.
\]

\newpage

{\Large\bfseries\color{ink}4. Use the vertex reports to fix the gauge signs}\par\vspace{3mm}

The upper wires below are the already measured bits, not surviving
vertex qubits. A double-wire control applies Z when its bit is one.
\par\noindent\begin{minipage}[t]{.485\linewidth}\centering\textbf{$b_b$ controls $Z_{A_7}$}

\begin{center}\begin{adjustbox}{max width=\linewidth}
\begin{quantikz}[row sep={.65cm,between origins},column sep=.65cm]
\lstick{$\color{groupb}b_b$} & \cwbend{1} & \cw \\
\lstick{$\color{gauge}A_{7}$} & \gate{Z} & \qw
\end{quantikz}
\end{adjustbox}\end{center}
\end{minipage}\hfill\begin{minipage}[t]{.485\linewidth}\centering\textbf{$b_c$ controls $Z_{A_{11}}$}

\begin{center}\begin{adjustbox}{max width=\linewidth}
\begin{quantikz}[row sep={.65cm,between origins},column sep=.65cm]
\lstick{$\color{groupc}b_c$} & \cwbend{1} & \cw \\
\lstick{$\color{gauge}A_{11}$} & \gate{Z} & \qw
\end{quantikz}
\end{adjustbox}\end{center}
\end{minipage}\par

\vspace{8mm}
\par\noindent\begin{minipage}[t]{.485\linewidth}\centering\textbf{$b_d$ controls $Z_{A_3}$}

\begin{center}\begin{adjustbox}{max width=\linewidth}
\begin{quantikz}[row sep={.65cm,between origins},column sep=.65cm]
\lstick{$\color{groupd}b_d$} & \cwbend{1} & \cw \\
\lstick{$\color{gauge}A_{3}$} & \gate{Z} & \qw
\end{quantikz}
\end{adjustbox}\end{center}
\end{minipage}\hfill\begin{minipage}[t]{.485\linewidth}\centering\textbf{$b_e$ controls $Z_{A_0}$}

\begin{center}\begin{adjustbox}{max width=\linewidth}
\begin{quantikz}[row sep={.65cm,between origins},column sep=.65cm]
\lstick{$\color{groupe}b_e$} & \cwbend{1} & \cw \\
\lstick{$\color{gauge}A_{0}$} & \gate{Z} & \qw
\end{quantikz}
\end{adjustbox}\end{center}
\end{minipage}\par

\[
F_b=Z_{A_7}^{b_b}Z_{A_{11}}^{b_c}Z_{A_3}^{b_d}Z_{A_0}^{b_e}.
\]
Apply these four possible gates in the displayed order. There is no
root-a correction. The already imposed even-parity condition fixes its
remaining Gauss sign.
\subhead{Release the five vertex qubits}

\begin{center}\begin{adjustbox}{max width=\linewidth}
\begin{quantikz}[row sep={.5cm,between origins},column sep=.65cm]
\lstick{$\color{groupa}v$} & \gate{\lvert0\rangle} & \rstick{$\text{helper pool}$}\qw
\end{quantikz}
\end{adjustbox}\end{center}

After their readout, all five vertices may be reset for helper reuse.
From this point on, a--e name graph locations and data groups; they are
not five live gauge-check ancillas.
\[
W=\frac14\sum_{t:t_a=0}|\delta t\rangle_A V_t,
\qquad |\text{gauged}\rangle\ \propto\ W\frac{I+U}{2}|\psi\rangle.
\]

\newpage

{\Large\bfseries\color{ink}5a. While gauged: measure the paired code checks}\par\vspace{3mm}


\[
N_0^X=X_{B_0}X_{B_1}X_{B_2}X_{B_3}
X_{C_0}X_{C_1}X_{C_2}X_{C_3}.
\]
Prepare \textbf{one eight-leaf cat} across both panels. These couplings
produce one eight-leaf parity, not two four-leaf results.
\par\noindent\begin{minipage}[t]{.485\linewidth}\centering\textbf{Leaves 0--3: B targets}

\begin{center}\begin{adjustbox}{max width=\linewidth}
\begin{quantikz}[row sep={.46cm,between origins},column sep=.22cm]
\lstick{$\color{cat}\ell_{0}$} & \ctrl{1} & \qw & \qw & \qw & \gate{H} & \meter{} & \rstick{$\lambda_{0}$}\cw \\
\lstick{$\color{groupa}B_0$} & \targ{} & \qw & \qw & \qw & \qw & \qw & \qw \\
\lstick{$\color{cat}\ell_{1}$} & \qw & \ctrl{1} & \qw & \qw & \gate{H} & \meter{} & \rstick{$\lambda_{1}$}\cw \\
\lstick{$\color{groupa}B_1$} & \qw & \targ{} & \qw & \qw & \qw & \qw & \qw \\
\lstick{$\color{cat}\ell_{2}$} & \qw & \qw & \ctrl{1} & \qw & \gate{H} & \meter{} & \rstick{$\lambda_{2}$}\cw \\
\lstick{$\color{groupa}B_2$} & \qw & \qw & \targ{} & \qw & \qw & \qw & \qw \\
\lstick{$\color{cat}\ell_{3}$} & \qw & \qw & \qw & \ctrl{1} & \gate{H} & \meter{} & \rstick{$\lambda_{3}$}\cw \\
\lstick{$\color{groupa}B_3$} & \qw & \qw & \qw & \targ{} & \qw & \qw & \qw
\end{quantikz}
\end{adjustbox}\end{center}
\end{minipage}\hfill\begin{minipage}[t]{.485\linewidth}\centering\textbf{Leaves 4--7: C targets}

\begin{center}\begin{adjustbox}{max width=\linewidth}
\begin{quantikz}[row sep={.46cm,between origins},column sep=.22cm]
\lstick{$\color{cat}\ell_{4}$} & \ctrl{1} & \qw & \qw & \qw & \gate{H} & \meter{} & \rstick{$\lambda_{4}$}\cw \\
\lstick{$\color{groupa}C_0$} & \targ{} & \qw & \qw & \qw & \qw & \qw & \qw \\
\lstick{$\color{cat}\ell_{5}$} & \qw & \ctrl{1} & \qw & \qw & \gate{H} & \meter{} & \rstick{$\lambda_{5}$}\cw \\
\lstick{$\color{groupa}C_1$} & \qw & \targ{} & \qw & \qw & \qw & \qw & \qw \\
\lstick{$\color{cat}\ell_{6}$} & \qw & \qw & \ctrl{1} & \qw & \gate{H} & \meter{} & \rstick{$\lambda_{6}$}\cw \\
\lstick{$\color{groupa}C_2$} & \qw & \qw & \targ{} & \qw & \qw & \qw & \qw \\
\lstick{$\color{cat}\ell_{7}$} & \qw & \qw & \qw & \ctrl{1} & \gate{H} & \meter{} & \rstick{$\lambda_{7}$}\cw \\
\lstick{$\color{groupa}C_3$} & \qw & \qw & \qw & \targ{} & \qw & \qw & \qw
\end{quantikz}
\end{adjustbox}\end{center}
\end{minipage}\par

\[
y(N_0^X)=\bigoplus_{j=0}^{7}\lambda_j;\qquad \text{require }y(N_0^X)=0.
\]
\subhead{The complete family}
For $R_s=(q_0<\cdots<q_{h-1})$, use $2h$ leaves, with ports
\[
\ell_j\to B_{q_j},\qquad \ell_{h+j}\to C_{q_j},\qquad j=0,\ldots,h-1.
\]
For $N_s^X$ use CNOT; for $N_s^Z$ replace each CNOT target by a CZ dot.
Measure $N_0^X,\ldots,N_8^X$, then $N_0^Z,\ldots,N_8^Z$.
The largest such cat has sixteen leaves.
\medskip

Do not measure $S_{B,s}$ and $S_{C,s}$ separately while gauged. The
required observable is their product. Cycle and Gauss checks follow,
then the entire 39-check round is repeated, three complete rounds total.

\newpage

{\Large\bfseries\color{ink}5b. While gauged: measure the sixteen cycle checks}\par\vspace{3mm}

Use one fresh verified cat for each cycle. All target gates are CZs.
The measurement acts only on gauge qubits.
\par\noindent\begin{minipage}[t]{.485\linewidth}\centering\textbf{$L_{(7,15)}=Z_{A_7}Z_{A_{15}}$}

\begin{center}\begin{adjustbox}{max width=\linewidth}
\begin{quantikz}[row sep={.46cm,between origins},column sep=.22cm]
\lstick{$\color{cat}\ell_{0}$} & \ctrl{1} & \qw & \gate{H} & \meter{} & \rstick{$\lambda_{0}$}\cw \\
\lstick{$\color{gauge}A_7$} & \control{} & \qw & \qw & \qw & \qw \\
\lstick{$\color{cat}\ell_{1}$} & \qw & \ctrl{1} & \gate{H} & \meter{} & \rstick{$\lambda_{1}$}\cw \\
\lstick{$\color{gauge}A_{15}$} & \qw & \control{} & \qw & \qw & \qw
\end{quantikz}
\end{adjustbox}\end{center}
\end{minipage}\hfill\begin{minipage}[t]{.485\linewidth}\centering\textbf{$L_{(7,8,11)}=Z_{A_7}Z_{A_8}Z_{A_{11}}$}

\begin{center}\begin{adjustbox}{max width=\linewidth}
\begin{quantikz}[row sep={.46cm,between origins},column sep=.22cm]
\lstick{$\color{cat}\ell_{0}$} & \ctrl{1} & \qw & \qw & \gate{H} & \meter{} & \rstick{$\lambda_{0}$}\cw \\
\lstick{$\color{gauge}A_7$} & \control{} & \qw & \qw & \qw & \qw & \qw \\
\lstick{$\color{cat}\ell_{1}$} & \qw & \ctrl{1} & \qw & \gate{H} & \meter{} & \rstick{$\lambda_{1}$}\cw \\
\lstick{$\color{gauge}A_8$} & \qw & \control{} & \qw & \qw & \qw & \qw \\
\lstick{$\color{cat}\ell_{2}$} & \qw & \qw & \ctrl{1} & \gate{H} & \meter{} & \rstick{$\lambda_{2}$}\cw \\
\lstick{$\color{gauge}A_{11}$} & \qw & \qw & \control{} & \qw & \qw & \qw
\end{quantikz}
\end{adjustbox}\end{center}
\end{minipage}\par

\subhead{All sixteen supports, in measurement order}
\[
\begin{aligned}
&(7,15),\ (11,12),\ (3,6),\ (0,10),\ (8,19),\\
&(14,16),\ (4,9),\ (2,13),\ (1,18),\ (5,17),\\
&(7,8,11),\ (7,14,3),\ (7,4,0),\\
&(11,2,3),\ (11,1,0),\ (3,5,0).
\end{aligned}
\]
Within a check, couple ports in increasing edge-label order, then measure
the leaves. Require the single parity $y(L_\eta)=\bigoplus_j\lambda_j=0$.
The six triangles together with the ten parallel-edge pairs form a
cycle basis. After these sixteen checks, measure the five Gauss checks.

\newpage

{\Large\bfseries\color{ink}5c. While gauged: the full local Gauss-check circuit}\par\vspace{3mm}


\[
G_a=(X_{A_0}X_{A_3}X_{A_6}X_{A_7}X_{A_{10}}X_{A_{11}}X_{A_{12}}X_{A_{15}})T_a.
\]
Prepare \textbf{one twelve-leaf cat} before all four panels. Eight ports
are CNOTs to gauge qubits; four are Fredkins to data pairs.
\par\noindent\begin{minipage}[t]{.485\linewidth}\centering\textbf{Gauge ports 0--3}

\begin{center}\begin{adjustbox}{max width=\linewidth}
\begin{quantikz}[row sep={.46cm,between origins},column sep=.22cm]
\lstick{$\color{cat}\ell_{0}$} & \ctrl{1} & \qw & \qw & \qw & \gate{H} & \meter{} & \rstick{$\lambda_{0}$}\cw \\
\lstick{$\color{gauge}A_{0}$} & \targ{} & \qw & \qw & \qw & \qw & \qw & \qw \\
\lstick{$\color{cat}\ell_{1}$} & \qw & \ctrl{1} & \qw & \qw & \gate{H} & \meter{} & \rstick{$\lambda_{1}$}\cw \\
\lstick{$\color{gauge}A_{3}$} & \qw & \targ{} & \qw & \qw & \qw & \qw & \qw \\
\lstick{$\color{cat}\ell_{2}$} & \qw & \qw & \ctrl{1} & \qw & \gate{H} & \meter{} & \rstick{$\lambda_{2}$}\cw \\
\lstick{$\color{gauge}A_{6}$} & \qw & \qw & \targ{} & \qw & \qw & \qw & \qw \\
\lstick{$\color{cat}\ell_{3}$} & \qw & \qw & \qw & \ctrl{1} & \gate{H} & \meter{} & \rstick{$\lambda_{3}$}\cw \\
\lstick{$\color{gauge}A_{7}$} & \qw & \qw & \qw & \targ{} & \qw & \qw & \qw
\end{quantikz}
\end{adjustbox}\end{center}
\end{minipage}\hfill\begin{minipage}[t]{.485\linewidth}\centering\textbf{Gauge ports 4--7}

\begin{center}\begin{adjustbox}{max width=\linewidth}
\begin{quantikz}[row sep={.46cm,between origins},column sep=.22cm]
\lstick{$\color{cat}\ell_{4}$} & \ctrl{1} & \qw & \qw & \qw & \gate{H} & \meter{} & \rstick{$\lambda_{4}$}\cw \\
\lstick{$\color{gauge}A_{10}$} & \targ{} & \qw & \qw & \qw & \qw & \qw & \qw \\
\lstick{$\color{cat}\ell_{5}$} & \qw & \ctrl{1} & \qw & \qw & \gate{H} & \meter{} & \rstick{$\lambda_{5}$}\cw \\
\lstick{$\color{gauge}A_{11}$} & \qw & \targ{} & \qw & \qw & \qw & \qw & \qw \\
\lstick{$\color{cat}\ell_{6}$} & \qw & \qw & \ctrl{1} & \qw & \gate{H} & \meter{} & \rstick{$\lambda_{6}$}\cw \\
\lstick{$\color{gauge}A_{12}$} & \qw & \qw & \targ{} & \qw & \qw & \qw & \qw \\
\lstick{$\color{cat}\ell_{7}$} & \qw & \qw & \qw & \ctrl{1} & \gate{H} & \meter{} & \rstick{$\lambda_{7}$}\cw \\
\lstick{$\color{gauge}A_{15}$} & \qw & \qw & \qw & \targ{} & \qw & \qw & \qw
\end{quantikz}
\end{adjustbox}\end{center}
\end{minipage}\par

\vspace{3mm}
\par\noindent\begin{minipage}[t]{.485\linewidth}\centering\textbf{Data-pair ports 8--9}

\begin{center}\begin{adjustbox}{max width=\linewidth}
\begin{quantikz}[row sep={.46cm,between origins},column sep=.24cm]
\lstick{$\color{cat}\ell_{8}$} & \ctrl{1} & \qw & \gate{H} & \meter{} & \rstick{$\lambda_{8}$}\cw \\
\lstick{$\color{groupa}B_{0}$} & \swap{1} & \qw & \qw & \qw & \qw \\
\lstick{$\color{groupa}C_{0}$} & \targX{} & \qw & \qw & \qw & \qw \\
\lstick{$\color{cat}\ell_{9}$} & \qw & \ctrl{1} & \gate{H} & \meter{} & \rstick{$\lambda_{9}$}\cw \\
\lstick{$\color{groupa}B_{1}$} & \qw & \swap{1} & \qw & \qw & \qw \\
\lstick{$\color{groupa}C_{1}$} & \qw & \targX{} & \qw & \qw & \qw
\end{quantikz}
\end{adjustbox}\end{center}
\end{minipage}\hfill\begin{minipage}[t]{.485\linewidth}\centering\textbf{Data-pair ports 10--11}

\begin{center}\begin{adjustbox}{max width=\linewidth}
\begin{quantikz}[row sep={.46cm,between origins},column sep=.24cm]
\lstick{$\color{cat}\ell_{10}$} & \ctrl{1} & \qw & \gate{H} & \meter{} & \rstick{$\lambda_{10}$}\cw \\
\lstick{$\color{groupa}B_{2}$} & \swap{1} & \qw & \qw & \qw & \qw \\
\lstick{$\color{groupa}C_{2}$} & \targX{} & \qw & \qw & \qw & \qw \\
\lstick{$\color{cat}\ell_{11}$} & \qw & \ctrl{1} & \gate{H} & \meter{} & \rstick{$\lambda_{11}$}\cw \\
\lstick{$\color{groupa}B_{3}$} & \qw & \swap{1} & \qw & \qw & \qw \\
\lstick{$\color{groupa}C_{3}$} & \qw & \targX{} & \qw & \qw & \qw
\end{quantikz}
\end{adjustbox}\end{center}
\end{minipage}\par

\[
y(G_a)=\bigoplus_{j=0}^{11}\lambda_j;\qquad \text{require }y(G_a)=0.
\]
The panels are port fragments, not four check measurements. Execute all
twelve controlled factors before the leaf readout. For v=b,c,d,e replace
$(E_a,J_a)$ by $(E_v,J_v)$, retaining the same twelve-port pattern.
\[
G_v=X_{A,E_v}T_v,\qquad G_aG_bG_cG_dG_e=U.
\]
One full round is $18\ N$ checks, $16\ L$ checks, then $5\ G$ checks.
Repeat that entire round three times, using fresh cats throughout.

\newpage

{\Large\bfseries\color{ink}6. Ungauge: read all edges, decode, and correct}\par\vspace{3mm}

Measure every $A_q$ in Z, increasing $q=0,\ldots,19$. Collect all
twenty reported bits before choosing any SWAP correction.

\begin{center}\begin{adjustbox}{max width=\linewidth}
\begin{quantikz}[row sep={.6cm,between origins},column sep=.7cm]
\lstick{$\color{gauge}A_q$} & \meter{} & \rstick{$m_q$}\cw
\end{quantikz}
\end{adjustbox}\end{center}

\begin{center}\begin{adjustbox}{max width=\linewidth}
\begin{quantikz}[row sep={.28cm,between origins},column sep=.8cm]
\lstick{$m_0,\ldots,m_{19}$} & \gate[cwires={1}]{\operatorname{Dec}_{\mathcal W}} & \rstick{$\widehat t_a,\ldots,\widehat t_e$}\cw
\end{quantikz}
\end{adjustbox}\end{center}

The double-wire box is classical computation. Enumerate all sixteen cut
words $w=\delta t$, with $t_a=0$. Choose the word $\widehat w$ within
Hamming distance two of m, or reject if none exists.
\[
\widehat t_a=0,\quad\widehat t_b=\widehat w_7,\quad
\widehat t_c=\widehat w_{11},\quad\widehat t_d=\widehat w_3,\quad
\widehat t_e=\widehat w_0.
\]
\subhead{A classical bit controls each ordinary data SWAP}

\begin{center}\begin{adjustbox}{max width=\linewidth}
\begin{quantikz}[row sep={.65cm,between origins},column sep=.8cm]
\lstick{$\widehat t_v$} & \cwbend{1} & \cw \\
\lstick{$\color{groupa}B_q$} & \swap{1} & \qw \\
\lstick{$\color{groupa}C_q$} & \targX{} & \qw
\end{quantikz}
\end{adjustbox}\end{center}

For each $q\in J_v$, apply SWAP$(B_q,C_q)$ when $\widehat t_v=1$,
omitting it otherwise. Execute $q=0,\ldots,19$ in order. These are not
coherent Fredkins: their controls are classical bits. Discard/reset A.
\[
\text{Decode all twenty reports first; do not use the four raw tree bits.}
\]

\newpage

{\Large\bfseries\color{ink}7--8. Filter the original code, then recover}\par\vspace{3mm}

Reuse the four-body Pauli circuit on page 2, extended to eight
targets for the upper checks. A fresh verified cat is required each time.
\[
\mathcal O=(S^X_{B,0:8},S^Z_{B,0:8},S^X_{C,0:8},S^Z_{C,0:8}).
\]
\subhead{7. Terminal positive filter}
Measure the complete 36-check family $\mathcal O$ in three rounds.
Require all $3\times36$ check bits zero. A nonzero syndrome rejects.
\subhead{8. Subsequent syndrome recovery}
Measure $\mathcal O$ repeatedly, retaining either check sign. Require
cat-verification bits zero. Apply no correction between rounds.
\[
\tau=\min\{r\in\{3,\ldots,9\}:\mathbf s^{(r)}=
\mathbf s^{(r-1)}=\mathbf s^{(r-2)}\}.
\]
Reject if there is no identical triple by round nine. For each block D,
choose a minimum-weight Pauli $R_D$, of weight at most two, matching its
eighteen stopping-syndrome bits; reject if none exists. Use the fixed
lookup in \texttt{c4\_helix\_intrinsic\_decoder.py}.
\subhead{The physical correction: one-qubit gates only}

\begin{center}\begin{adjustbox}{max width=\linewidth}
\begin{quantikz}[row sep={.65cm,between origins},column sep=.8cm]
\lstick{$x_{D,q},z_{D,q}$} & \cw\vcw{1} & \cw \\
\lstick{$\color{groupa}D_q$} & \gate{P_{x,z}} & \qw
\end{quantikz}
\end{adjustbox}\end{center}

The classical pair selects $P_{0,0}=I$, $P_{1,0}=X$, $P_{0,1}=Z$,
or $P_{1,1}=Y$. Apply the selected single-qubit Pauli, visiting physical
labels in order, B then C; omit identity factors. For Y use one Y gate.
The gate acts on one data qubit; the syndrome lookup is classical.

\newpage

{\Large\bfseries\color{ink}9. Decode the final B readout and retain magic in C}\par\vspace{3mm}

For $q=0,\ldots,19$, measure B in Z. C remains a quantum output.
The horizontal C wire also represents its counted waiting idles.

\begin{center}\begin{adjustbox}{max width=\linewidth}
\begin{quantikz}[row sep={.65cm,between origins},column sep=.7cm]
\lstick{$\color{groupa}B_q$} & \meter{} & \rstick{$n_q$}\cw \\
\lstick{$\color{groupa}C_q$} & \qw & \rstick{$\text{retain}$}\qw
\end{quantikz}
\end{adjustbox}\end{center}

\begin{center}\begin{adjustbox}{max width=\linewidth}
\begin{quantikz}[row sep={.28cm,between origins},column sep=.6cm]
\lstick{$n_0,\ldots,n_{19}$} & \gate[cwires={1}]{\operatorname{Dec}_{B}} & \gate[cwires={1}]{\ell(n\oplus e)} & \rstick{$\text{keep iff }0$}\cw
\end{quantikz}
\end{adjustbox}\end{center}

\[
\begin{aligned}
h_s(n)&=\bigoplus_{q\in R_s}n_q,\qquad s=0,\ldots,8,\\
&\text{choose minimum-weight }e,\quad |e|\le2,\quad h(e)=h(n),\\
n'&=n\oplus e,\\
\ell(n')&=n'_0\oplus n'_3\oplus n'_6\oplus n'_7\oplus n'_{18}\oplus n'_{19}.
\end{aligned}
\]
Reject if the correction e does not exist. Fix readout ties by
lexicographic order of the flipped-position lists. Keep only
$\ell(n')=0$. Correction of n is classical; no postmeasurement X gates
are applied to B. Discard B and retain the twenty-qubit C block.
\subhead{Ideal logical output}
\[
\frac{2|00\rangle+|01\rangle+|10\rangle}{\sqrt6}
\ \xrightarrow{\ \overline Z^B_0=+1\ }\ 
|0\rangle_B\frac{2|0\rangle+|1\rangle}{\sqrt5}_C.
\]
Both second logical qubits are spectators in zero.
The retained output is $|\overline{M0}\rangle_C$, up to the physical
correctable errors allowed by the linked two-fault analysis.

\newpage

{\Large\bfseries\color{ink}Reference: supports, order, and resources}\par\vspace{3mm}


Each support defines $S^X_{D,s}=\prod_{q\in R_s}X_{D_q}$ and
$S^Z_{D,s}=\prod_{q\in R_s}Z_{D_q}$. Measure only $s=0,\ldots,8$
as independent code checks; use row 9 for graph incidence.
\begin{center}\begin{tabular}{cl}\toprule
s & $R_s$\\\midrule
0 & 0, 1, 2, 3 \\
1 & 4, 5, 6, 7 \\
2 & 8, 9, 10, 11 \\
3 & 12, 13, 14, 15 \\
4 & 16, 17, 18, 19 \\
5 & 0, 3, 6, 7, 10, 11, 12, 15 \\
6 & 4, 7, 8, 9, 14, 15, 16, 19 \\
7 & 1, 2, 8, 11, 12, 13, 18, 19 \\
8 & 2, 3, 5, 6, 13, 14, 16, 17 \\
9 & 0, 1, 4, 5, 9, 10, 17, 18 \\
\bottomrule\end{tabular}\end{center}
\[
\begin{aligned}
\overline X_0&=X_{\{2,3,8,11,12,15\}},\\
\overline Z_0&=Z_{\{0,3,6,7,18,19\}},\qquad
\overline Z_1=Z_{\{2,3,8,11,12,15\}}.
\end{aligned}
\]
\subhead{One attempt, in execution order}
\begin{enumerate}\setlength\itemsep{1pt}
\item Prepare B, then C: three positive 20-check rounds each.
\item Prepare graph; three positive 25-check resource rounds.
\item Twenty enrichment Fredkins; measure the five vertices in X.
\item Require even vertex parity; apply the four possible gauge-Z gates.
\item Three positive gauged rounds: $18N$, $16L$, $5G$, in that order.
\item Gauge-Z readout; decode; apply ordinary data SWAPs.
\item Three positive 36-check original-code rounds.
\item Recovery: first identical syndrome triple, at most nine rounds.
\item Complete B-Z readout; decode; keep logical zero and retain C.
\end{enumerate}
\textbf{Allocation:} 40 data + 20 gauge + 17 reusable helpers = 77 qubits.
Five helpers initially serve as vertices. The largest cat has 16 leaves
plus one verifier. \textbf{Fredkins:} $20+3\times5\times4=80$.
\medskip

All gates and measurements are scheduled serially; paired panels are
fragments of a circuit, not additional parallel hardware. Unused live
qubits have counted idles. Any required failed test rejects the whole
attempt. The global fault allowance is two faulty locations total,
including helpers, reports, recovery, and readout, as specified in the
companion analysis. These drawings do not change that protocol or prove
a new fault-tolerance result.


\end{document}
